\chapter{Data Collection Methods and Sources}


\section{Data Sources}
\textbf{Bus Map Data}: The data is collected by scraping directly from the \url{https://map.busmap.vn/hcm} web.


\textbf{Traffic Camera Data}: The data is collected by scraping directly from the \url{https://giaothong.hochiminhcity.gov.vn} web.

\section{Data Collection Methods}
\subsection{Bus Map Data}
Step-by-step implementation process:
\begin{itemize}
\item \textbf{Raw Capture (Interception):} A script uses playwright to launch a browser context and listen for network responses containing "api" or "route". These responses are captured into a file named api\_responses.json.
\item \textbf{Decryption:} The API responses are encrypted. The system uses the AES-CBC algorithm (via Crypto.Cipher.AES) to decrypt the hexadecimal payloads. The decryption key is extracted dynamically from the initial data packet.
\item \textbf{Static Resource Download:} The system extracts URLs for static resources (routes, route info, stations) from the decrypted data and downloads them to a raw directory. The extractor automatically handles various compression formats, including Gzip and Zlib.
\item \textbf{Tools used:} Python, Playwright (automation), PyCryptodome (decryption), Requests (HTTP calls).
\item \textbf{Storage format:} The raw data is stored as JSON or binary .xyz files, while the final output is stored as CSV.
\end{itemize}

\begin{figure}[h!]
\centering

\begin{tikzpicture}[
  node distance = 18mm,
  every node/.style={font=\small},
  proc/.style = {
      rectangle,
      text width=0.5\linewidth,  
      minimum height=11mm,
      text centered,
      draw=black,
      very thick,
      rounded corners=3pt,
      font=\small,
  },
  arr/.style = {-{Stealth[length=3mm]}, thick},
  note/.style = {
      rectangle,
      rounded corners,
      draw=black!60,
      fill=black!5,
      inner sep=6pt,
      font=\footnotesize,
      align=left,
      text width=0.35\linewidth,
  }
]

% ==== NODES (VERTICAL FLOW WITH A4 WIDTH) ====
\node[proc, fill=blue!20] (capture) {%
  \textbf{Raw Capture (Interception)}\newline
  Playwright captures network responses ("api", "route")\newline
  Stored as \texttt{api\_responses.json}%
};

\node[proc, fill=yellow!25, below=18mm of capture] (decrypt) {%
  \textbf{Decryption}\newline
  AES-CBC using PyCryptodome\newline
  Dynamic key extraction from initial data packet%
};

\node[proc, fill=green!20, below=18mm of decrypt] (extract) {%
  \textbf{Static Resource Download}\newline
  Extract route/station URLs\newline
  Auto-handle gzip/zlib decompression%
};

\node[proc, fill=red!20, below=18mm of extract] (store) {%
  \textbf{Storage \& Final Output}\newline
  Raw: JSON / binary .xyz\newline
  Final: CSV%
};

% ==== ARROWS ====
\draw[arr] (capture) -- (decrypt);
\draw[arr] (decrypt) -- (extract);
\draw[arr] (extract) -- (store);

% ==== NOTES ====
\node[note, right=15mm of decrypt] (tools) {%
  \textbf{Tools Used:}\newline
  Python\newline Playwright (automation)\newline PyCryptodome (AES-CBC)\newline Requests (HTTP)%
};

\node[note, right=15mm of store] (formats) {%
  \textbf{Storage Formats:}\newline
  Raw: JSON, \texttt{.xyz}\newline
  Final: CSV%
};

\end{tikzpicture}

\caption{Diagram for Bus Map Data Collection}
\end{figure}


\subsection{Traffic Camera Data}
Step-by-step implementation process:
\begin{itemize}
\item \textbf{Command-Line Argument Parsing and Initialization:}
\begin{itemize}
\item  Use argparse to parse arguments: --cameras-file (JSON with camera list), --interval (cycle seconds, default 120), --max-workers (threads, default 6), --output-dir (storage root, default ./camera\_frames). 
\item Load camera entries from JSON file. 
\item Optionally load camera\_locations.json for GPS and location metadata enrichment.
\end{itemize}

\item \textbf{Multi-Camera Manager Setup:}
\begin{itemize}
    \item Instantiate MultiCameraLoop with camera list, interval, workers, output directory, and locations.
    \item Create an HCMCCameraCrawler instance per camera, each with a persistent requests.Session, custom headers (mimicking browser), timeout (60s), and deduplication/stats tracking.
\end{itemize}

\item \textbf{Signal Handling for Graceful Shutdown:}
Register handlers for SIGINT/SIGTERM to set a stop event, allowing completion of the current cycle.

\item \textbf{Continuous Looping Process:}
\begin{itemize}
\item Enter an infinite loop until stop event is set.
\item For each cycle:
\begin{itemize}
\item Print cycle start timestamp (UTC+7).
\item Use ThreadPoolExecutor (max\_workers threads) to concurrently execute crawl\_once for each camera.
\item In crawl\_once:
\begin{itemize}
\item Increment attempt counter.
\item Fetch image: Construct URL with camera ID and millisecond timestamp; send GET request; validate response (status 200, image content-type, size >1000 bytes).
\item If successful: Compute SHA256 hash; check for duplication (matches previous hash); increment stats.
\item If unique: Generate timestamped filename; save JPG image and JSON metadata (camera details, timestamps, hash, size, GPS if available).
\item Cleanup: Sort images by modification time; delete all but the newest 3 (configurable) per camera folder, including companion JSON.
\end{itemize}
\item Wait for all tasks; aggregate and print cycle stats (attempts, successes, failures, saves, duplicates).
\item Sleep for remaining interval time (in 1-second increments to check stop event).
\end{itemize}   
\end{itemize}

\item \textbf{Shutdown and Final Reporting:}
On stop: Complete cycle, print final per-camera stats (including uptime), and exit.
\end{itemize}

\begin{figure}[!h]
\centering

\begin{tikzpicture}[
  node distance = 18mm,
  every node/.style={font=\small},
  proc/.style = {
      rectangle,
      text width=0.6\linewidth,
      minimum height=12mm,
      text centered,
      draw=black,
      very thick,
      rounded corners=3pt,
      font=\small,
      align=left,
      inner sep=4pt,
  },
  arr/.style = {-{Stealth[length=3mm]}, thick},
  note/.style = {
      rectangle,
      rounded corners,
      draw=black!60,
      fill=black!5,
      inner sep=6pt,
      font=\footnotesize,
      align=left,
      text width=0.35\linewidth,
  }
]

% ==== NODES (VERTICAL FLOW) ====
\node[proc, fill=blue!20] (argparse) {%
  \textbf{Command-Line Argument Parsing \& Initialization}\newline
  - argparse: --cameras-file, --interval, --max-workers, --output-dir\newline
  - Load camera entries from JSON\newline
  - Optional: load camera\_locations.json%
};

\node[proc, fill=yellow!25, below=18mm of argparse] (manager) {%
  \textbf{Multi-Camera Manager Setup}\newline
  - MultiCameraLoop: cameras, interval, workers, output dir, locations\newline
  - HCMCCameraCrawler per camera: persistent session, headers, timeout, dedup/stats%
};

\node[proc, fill=green!20, below=18mm of manager] (signal) {%
  \textbf{Signal Handling for Graceful Shutdown}\newline
  - Register SIGINT/SIGTERM handlers to set stop event%
};

\node[proc, fill=orange!20, below=18mm of signal] (loop) {%
\textbf{Continuous Looping Process}\newline
  - Infinite loop until stop event\newline
  - For each cycle:\newline
    \quad - Print start timestamp (UTC+7)\newline
    \quad - ThreadPoolExecutor for concurrent crawl\_once per camera\newline
    \quad - In crawl\_once: fetch image, validate, deduplicate, save JPG + JSON, cleanup old files\newline
    \quad - Wait tasks, print cycle stats, sleep for remaining interval%
};

\node[proc, fill=red!20, below=18mm of loop] (shutdown) {%
  \textbf{Shutdown \& Final Reporting}\newline
  - Complete cycle, print per-camera stats (uptime), exit%
};

% ==== ARROWS ====
\draw[arr] (argparse) -- (manager);
\draw[arr] (manager) -- (signal);
\draw[arr] (signal) -- (loop);
\draw[arr] (loop) -- (shutdown);

% ==== NOTES ====
\node[note, right=15mm of argparse] (note1) {%
\textbf{Args Examples:}\newline
--cameras-file cameras.json\newline
--interval 120\newline
--max-workers 6\newline
--output-dir ./camera\_frames%
};

\node[note, right=15mm of loop] (note2) {%
\textbf{Image Handling:}\newline
- Validate: status 200, image type, size >1000B\newline
- Dedup via SHA256 hash\newline
- Save JPG + metadata JSON\newline
- Keep latest 3 images per camera%
};

\end{tikzpicture}

\caption{Diagram for Multi-Camera Crawling}
\end{figure}



% \section{Architecture}

% \subsection{System Components}

% \textbf{Scheduler:}
% \begin{itemize}
%     \item \textbf{Implementation}: Python-based task scheduling.
%     \item \textbf{Scheduling Strategy}: Configurable intervals (15 minutes) for camera frame collection.
%     \item \textbf{Task Queue}: Manages concurrent camera crawling jobs to prevent overload.
%     \item \textbf{Monitoring}: Tracks crawler execution status, logs failures, triggers alerts on consecutive failures.
% \end{itemize}

% \textbf{Crawling Engine:}
% \begin{itemize}
%     \item \textbf{Primary Tools}: Python Requests library for HTTP API calls, Playwright/Selenium for dynamic content.
%     %\item \textbf{Frontend Client}: React 17 application (4,981 lines) with axios v1.13.2 for API communication.
%     \item \textbf{Base URL}: \texttt{https://api.hcmus.fit} for all backend services.
%     \item \textbf{HTTP Client Configuration}: Axios interceptors for automatic token refresh on 401 responses, request/response logging, error handling.
% \end{itemize}


% \textbf{Error Handling and Retry Strategy:}
% \begin{itemize}
%     \item \textbf{Network Failures}: Exponential backoff retry (3 attempts with 2s, 4s, 8s delays).
%     \item \textbf{Authentication Errors}: \texttt{tokenManager.js} handles token expiration gracefully:
%     \begin{itemize}
%         \item \texttt{handleTokenExpiration()} clears \texttt{localStorage}, dispatches \texttt{auth-change} event.
%         \item Shows animated notification (gradient background, slideDown animation, 3s display).
%         \item Redirects to \texttt{/login} after 1s delay.
%         \item \texttt{isLoggingOut} mutex prevents concurrent logout calls.
%     \end{itemize}
%     \item \textbf{Camera Unavailability}: Skip failed camera, log error, continue with remaining cameras. Retry on next scheduled run.
%     \item \textbf{API Rate Limiting}: Respect \texttt{Retry-After} headers, implement request throttling (max N requests/second).
% \end{itemize}

% \section{Traffic Camera Crawling}

% \subsection{Collection Workflow}

% \textbf{Frame Capture Process:}
% \begin{enumerate}
%     \item Scheduler triggers camera collection job every 15 minutes.
%     \item For each of 702 cameras in \texttt{camera\_locations.json}:
%     \begin{itemize}
%         \item Fetch latest frame from camera API endpoint.
%         \item Validate image (non-zero size, valid JPEG header).
%         \item Save to \texttt{public/camera\_frames/\{camera\_id\}/}.
%     \end{itemize}
%     \item Generate metadata JSON with timestamp, hash (SHA-256), file size.
%     \item Update \texttt{camera\_frames} database table with new record.
%     \item Clean up old frames (optional: delete frames older than N days to manage storage).
% \end{enumerate}


% \textbf{Metadata Standardization:}
% \begin{itemize}
%     \item \textbf{Timestamps}: ISO 8601 format (\texttt{2025-11-16T14:34:43.593455Z}) in \texttt{timestamp\_iso} field, Unix epoch milliseconds in \texttt{timestamp} field.
%     \item \textbf{Camera ID}: Unique identifier from \texttt{camera\_locations.json}, used as primary key in \texttt{cameras} table and foreign key in \texttt{camera\_frames}.
%     \item \textbf{Hash}: SHA-256 digest for integrity verification and duplicate detection.
%     \item \textbf{Filename Convention}: \texttt{\{camera\_id\}\_\{date\}\_\{time\}\_\{hash\_prefix\}.jpg} for easy sorting and retrieval.
% \end{itemize}

% \textbf{Storage Optimization:}
% \begin{itemize}
%     \item Current storage: 35 MB for 702 frames (one per camera, latest only in local build).
%     \item Production storage: Estimated \textasciitilde70000 frames/day \texttimes{} 50 KB = \textasciitilde3.5 GB/day.
%     \item Time-series partitioning: Partition \texttt{camera\_frames} table by day or week.
%     \item Cold storage archival: Move frames older than 3 days to object storage with lifecycle policies.
% \end{itemize}

% \section{Web Crawling}

% \subsection{HTML Scraping}

% \textbf{Tools:}
% \begin{itemize}
%     \item \textbf{BeautifulSoup}: Python library for parsing HTML/XML documents.
%     \item \textbf{Playwright}: Headless browser automation for JavaScript-heavy sites.
%     \item \textbf{Selector Strategies}: CSS selectors, XPath for precise element targeting.
% \end{itemize}

% \textbf{Use Cases:}
% \begin{itemize}
%     \item Scraping bus schedules from \url{https://map.busmap.vn/hcm}
%     \item Collecting public transit stops, route descriptions, fare information.
% \end{itemize}

% \subsection{API Fetching (REST/JSON)}

% \textbf{Geocoding API:}
% \begin{itemize}
%     \item \textbf{Endpoint}: \texttt{POST /search} with JSON \texttt{\{address, lat, lon\}}.
%     \item \textbf{Provider}: Goong Maps API.
%     \item \textbf{Usage}: Forward geocoding (address to coordinates), reverse geocoding (coordinates to address).
%     \item \textbf{Response Caching}: Cache frequent queries (e.g., major landmarks) in Redis with 24-hour TTL.
% \end{itemize}

% \textbf{Routing API:}
% \begin{itemize}
%     \item \textbf{Endpoint}: \texttt{POST /route} with JSON \texttt{\{start, end, travelMode\}}.
%     \item \textbf{Travel Modes}: car, walk, bike.
%     \item \textbf{Response}: Route geometry, distance (meters), duration (seconds), step-by-step instructions.
%     \item \textbf{Integration}: \texttt{SearchBoxRoutes.js} component submits user queries, displays results on map.
% \end{itemize}

% \textbf{Bus Routing API:}
% \begin{itemize}
%     \item \textbf{Endpoint}: \texttt{GET /api/bus/route} with query parameters \texttt{start\_lat}, \texttt{start\_lng}, \texttt{end\_lat}, \texttt{end\_lng}, \texttt{max\_walk} (default 300m).
%     \item \textbf{Response}: Multi-leg itinerary with walking segments, bus transfers, route names, ETAs.
%     \item \textbf{Implementation}: \texttt{BusMapPage.js} with Material-UI search interface, route polyline visualization.
% \end{itemize}

% \subsection{Dynamic Content Scraping}

% \textbf{Headless Browser Techniques:}
% \begin{itemize}
%     \item \textbf{Playwright/Selenium}: For sites requiring JavaScript execution to render content.
%     \item \textbf{Wait Strategies}: Wait for specific elements (\texttt{page.waitForSelector}), network idle (\texttt{page.waitForLoadState('networkidle')}).
%     \item \textbf{Interaction Simulation}: Click buttons, fill forms, scroll to load infinite-scroll content.
% \end{itemize}

% \textbf{Example Scenario:}
% Scraping real-time bus locations from interactive transit map:
% \begin{enumerate}
%     \item Launch headless browser with Playwright.
%     \item Navigate to transit authority website.
%     \item Wait for map to load, extract bus marker coordinates from JavaScript variables or DOM.
%     \item Parse bus IDs, routes, current positions, timestamps.
%     \item Store in database for display in \texttt{BusMapPage}.
% \end{enumerate}

% \subsection{Anti-Ban Strategies and Request Throttling}

% \textbf{IP Rotation:}
% % \begin{itemize}
% %     \item Use proxy pools or VPN services to distribute requests across multiple IPs.
% %     \item Rotate user agents to mimic different browsers (Chrome, Firefox, Safari).
% % \end{itemize}

% \textbf{Request Throttling:}
% \begin{itemize}
%     \item Limit concurrent requests to max N per second (e.g., 5 req/s).
%     \item Introduce random delays between requests (1-3 seconds) to appear human-like.
%     \item Respect \texttt{robots.txt} directives and website terms of service.
% \end{itemize}

% \textbf{Session Management:}
% \begin{itemize}
%     \item Maintain cookies across requests to preserve session state.
%     \item Handle CSRF tokens if website requires them for form submissions.
% \end{itemize}

% \textbf{Error Detection:}
% \begin{itemize}
%     \item Monitor HTTP status codes: 429 (rate limit), 403 (forbidden), 503 (service unavailable).
%     \item Detect CAPTCHA challenges, pause crawler, send alert for manual intervention.
%     \item Log blocked IP addresses, rotate to new proxy on detection.
% \end{itemize}



% % \section{Summary}

% % The data crawling and collection pipeline employs a multi-layered architecture with robust error handling, validation, and cleaning mechanisms. Scheduler-driven camera frame collection, API-based data fetching, and web scraping techniques ensure comprehensive data acquisition. Anti-ban strategies and request throttling maintain crawler sustainability, while validation and deduplication steps guarantee data quality for downstream analytics and user-facing services.
